%% Cover letter for EMS submission — hydroflux methods paper.
%% Standalone document; compile with xelatex (matches the paper).
%%   xelatex cover_letter_ems
\documentclass[11pt]{article}

\usepackage{fontspec}
\setmainfont{Latin Modern Roman}
\setmonofont[Scale=0.95]{DejaVu Sans Mono}
\usepackage[a4paper, margin=2.5cm]{geometry}
\usepackage{parskip}
\usepackage{hyperref}
\usepackage{xcolor}
\hypersetup{colorlinks=true, linkcolor=black, urlcolor=blue!50!black,
            citecolor=blue!50!black}
\usepackage{enumitem}
\setlist[itemize]{leftmargin=*, itemsep=2pt, topsep=2pt}

\begin{document}

\noindent
\textbf{Cover letter} \hfill \today

\medskip
\noindent To: The Editors-in-Chief, \emph{Environmental Modelling \&
Software} (Prof.\ Min Chen; Prof.\ Sondoss El Sawah)

\medskip
Dear Editors,

I am submitting the manuscript \emph{``hydroflux: a well-balanced,
differentiable-by-design 2D shallow-water solver in Rust, verified
against analytical and community benchmarks and applied to a semiarid
Andean reach''} for consideration as a \textbf{Research Article} in
\emph{Environmental Modelling \& Software}.

The manuscript fits the journal's Aims \& Scope on three explicit axes:

\begin{enumerate}[leftmargin=*, label=\arabic*.]
\item \textbf{Environmental software with engineering rigor and
quantitative V\&V.} The solver is verified against a deliberate
hierarchy --- lake-at-rest to machine precision, the Thacker
oscillating paraboloid, the Stoker/Ritter dam-break, a radial
dam-break, steady Manning uniform flow, and the six UK Environment
Agency 2D benchmark tests --- and a matched head-to-head against
ANUGA on Stoker recovers the same accuracy class. A mesh-refinement
convergence study reports orders $L^1\ 1.81$ / $L^2\ 1.68$ on the
Thacker problem, consistent with the expected front-limited behaviour
of MUSCL schemes at moving wet/dry shorelines. All test inputs,
benchmark scripts, and verification outputs are reproducible from the
open-source repository.

\item \textbf{Real-world environmental application illustrating the
methodology.} The solver is applied to the 2017 Aluvión Atacama event
on the Río Huasco at Santa Juana --- a semiarid Andean reach with a
92-year DGA record --- on a 30~m DEM, using a spatially variable
Manning field derived from ESA WorldCover land cover. The result is
reported honestly as a one-day-peak sensitivity demonstration (not a
calibrated hindcast); the calibration companion is in preparation as
a separate submission.

\item \textbf{A generalizable software-engineering insight beyond the
specific application.} The solver is \emph{generic over the numeric
type} through a Rust trait abstraction, so that the identical code
path evaluates in \texttt{f64} for production and in forward-mode
dual numbers for gradient extraction. This zero-runtime-cost approach
to differentiability is, to our knowledge, the first such pattern
reported in a compiled environmental-modelling kernel; the lesson
transfers directly to any structured-grid hyperbolic solver in a
modern systems language. The well-balancedness, mass conservation
under wet/dry, and cell-mask early-skip optimisation are documented
with the discipline EMS readers expect, including a mass-conservation
discussion of why bounding-box skip-dry strategies break it (resolved
here by an inside-closure cell-mask).
\end{enumerate}

\textbf{Positioning relative to recent EMS contributions.} The
submission sits in a line of shallow-water solver papers EMS has
published recently (Rak et al.\ 2024, \emph{EMS} 177; Saleem \& Norman
2024, \emph{EMS} 180; Chen et al.\ 2025, \emph{EMS} 189; Li, Caviedes-
Voullième et al.\ 2024, \emph{EMS} 171). The differentiator is
\emph{not} a GPU speedup claim --- that is on the explicit roadmap
(§5) and not yet delivered. The differentiator is the combination of
(i) rigorous well-balanced numerics on a Rust+\texttt{Real}
foundation, (ii) a full verification hierarchy reproducible from the
public repository, and (iii) the differentiability pattern as a
forward-looking software-engineering contribution.

\textbf{Companion submission already at EMS.} The GeoTIFF I/O backbone
used by hydroflux (\texttt{surtgis}) is currently under Major Revision
at \emph{Environmental Modelling \& Software} (separate manuscript).
The two submissions are independent in claims and reviewable in
isolation; they form a coherent open-source stack, which we note for
the editor's awareness.

\textbf{Open source, licence, and reproducibility.} The solver is
released under the MIT OR Apache-2.0 dual licence at
\url{https://github.com/franciscoparrao/hydroflux}, with the benchmark
suite, the figure-generation scripts, and the Huasco application
notebook all included. A Zenodo snapshot DOI will be minted at
acceptance.

\textbf{Conflict of interest declaration.} The corresponding author is
the developer of the hydroflux solver. The head-to-head comparison
against ANUGA (§3.8) is reported with its outcome unaltered: ANUGA
achieves $L^1\ 2.6\%$ against the Ritter analytical at $\Delta x =
1$~m, hydroflux achieves $L^1\ 4.1\%$ at the same resolution; the gap
closes to $L^1\ 1.0\%$ under refinement to 400 cells. No result was
selected to favour hydroflux.

\textbf{Declaration on AI-assisted writing.} Generative AI tools were
used to assist drafting, refactoring, and reference verification,
with all content reviewed and validated by the authors. Three
references in an earlier draft (Hydrograd.jl, AegirJAX, r.avaflow~v4)
failed verify-refs / OpenAlex checks and were removed before submission.

The manuscript has not been submitted elsewhere and is not under
consideration by any other journal.

\textbf{Suggested reviewers} (no conflict of interest with the
authors):
\begin{itemize}
\item Prof.\ Daniel Caviedes-Voullième (Forschungszentrum Jülich) ---
GPU-accelerated shallow-water modelling, EMS author 2024.
\item Prof.\ Stephen Roberts (Australian National University) ---
ANUGA lead developer; relevant to the §3.8 head-to-head and the §5
roadmap.
\item Prof.\ Pilar García-Navarro (University of Zaragoza) --- Iber
solver lineage, well-balanced numerics on irregular topography.
\item Prof.\ Chaopeng Shen (Penn State) --- differentiable hydrology;
relevant to the differentiability framing of §1 and §5.
\item Prof.\ Cristián Escauriaza (Pontificia Universidad Católica de
Chile) --- Andean fluvial hydraulics; relevant to the Huasco
application context.
\end{itemize}

\textbf{Reviewers we would prefer to avoid}: none specifically. We
would only ask the editor to be mindful of reviewers with a direct
competing codebase under active development in Rust (a rare profile
at present).

We appreciate the editor's time and the reviewers' effort and look
forward to your decision.

\medskip
Sincerely,

\medskip
\noindent
\textbf{Francisco Parra} (corresponding author)\\
Postdoctoral researcher, Universidad de Santiago de Chile\\
\href{mailto:francisco.parra.o@usach.cl}{francisco.parra.o@usach.cl}
\quad·\quad ORCID 0009-0006-0435-1854

\medskip
\noindent
on behalf of co-authors\\
V.~Gil-Costa (UNSL--CONICET, San Luis, Argentina)\\
C.~Bonacic (Universidad de Santiago de Chile)\\
M.~Marín (Universidad de Santiago de Chile)

\end{document}
